\subsection{3. Thảo luận chung về các kết quả thực nghiệm}
\indent\indent Dựa trên các thực nghiệm được tiến hành từ mô hình đơn phương thức đến đa phương thức, kết hợp với các kiểm chứng trên các mô hình máy học truyền thống, các thảo luận tổng quát về hiệu quả của giải pháp đề xuất được tổng hợp như sau:
\begin{itemize}[label={--}, itemsep=1pt]
    \item \textbf{Hiệu quả của phương pháp đa phương thức:} Kết quả thực nghiệm cho thấy việc tích hợp thông tin văn bản và âm thanh thông qua kiến trúc MGAT giúp hệ thống đạt độ chính xác cao hơn so với các mô hình phân lớp chỉ dựa trên hình ảnh (CNN Baseline). Các đặc trưng ngữ nghĩa từ văn bản cùng với thông tin âm thanh đóng vai trò bổ trợ cho các đặc trưng hình ảnh, qua đó làm giảm hiện tượng nhập nhằng giữa các lớp có khung cảnh hình ảnh tương đồng và cải thiện độ tin cậy của dự đoán.
    
   \item \textbf{Chất lượng biểu diễn đặc trưng của MGAT:} Thông qua kết quả kiểm thử với các mô hình máy học truyền thống, trong đó Linear SVC đạt độ chính xác hơn 86\%, có thể thấy rằng kiến trúc MGAT học được các biểu diễn đặc trưng mang tính phân biệt cao. Các vector đặc trưng đầu ra hỗ trợ hiệu quả việc gom nhóm các mẫu dữ liệu cùng lớp trong không gian nhiều chiều, ngay cả trước khi áp dụng các bộ phân lớp phi tuyến phức tạp.

    \item \textbf{Sự tương thích giữa CNN backbone và bộ phân lớp:} Hai bộ phân lớp FastKAN và MLP cho thấy mức độ phụ thuộc cao vào đặc điểm của CNN backbone:
    \begin{itemize}[label={+}, itemsep=1pt]
        \item FastKAN đạt hiệu quả cao khi kết hợp với các CNN backbone có khả năng trích xuất đặc trưng sâu và giàu thông tin (chẳng hạn như DenseNet và ResNet). Khả năng xấp xỉ hàm linh hoạt của KAN cho phép khai thác hiệu quả không gian đặc trưng có độ biểu diễn cao.
        \item Ngược lại, MLP thể hiện sự ổn định và phù hợp hơn khi kết hợp với các CNN backbone nhẹ (như MobileNet Small và ShuffleNet). Trong trường hợp không gian đặc trưng đầu vào bị nén hoặc có mức độ thông tin hạn chế, cấu trúc đơn giản của MLP giúp giảm nguy cơ overfitting so với FastKAN.
    \end{itemize}

    \item \textbf{Hiệu quả của chiến lược suy luận phân tầng:} Sự khác biệt về hiệu năng giữa mô hình hai phương thức và mô hình ba phương thức cho thấy dữ liệu âm thanh đóng vai trò quan trọng nhưng chỉ thực sự hữu ích trong những ngữ cảnh nhất định. Chiến lược suy luận hai bước, trong đó bước đầu tiên thực hiện lọc nhãn thô nhằm tránh đưa thông tin âm thanh vào các lớp không phụ thuộc vào âm thanh, giúp duy trì độ chính xác ổn định khi triển khai thực tế, đồng thời đảm bảo sự nhất quán với kết quả đánh giá trên tập Test.
\end{itemize}

\indent Trong phạm vi của đề tài này, việc so sánh các mô hình chỉ được thực hiện trên khía cạnh không gian, cụ thể là số lượng tham số, mà chưa xem xét đến thời gian chạy. Lý do là mục tiêu chính của các thực nghiệm tập trung vào việc đánh giá hiệu quả sử dụng tham số của kiến trúc đề xuất. Thông qua việc áp dụng chiến lược học chuyển giao (Transfer Learning) bằng cách đóng băng các backbone, mô hình chỉ cần cập nhật một lượng tham số tương đối nhỏ (khoảng 4-5M) so với tổng lượng tham số mô hình (khoảng 28-40M). Cách tiếp cận này nhằm chứng minh rằng, chỉ với việc huấn luyện một phần mô hình có quy mô gọn nhẹ, hệ thống vẫn có thể xử lý hiệu quả tri thức từ các mạng tiền huấn luyện lớn để đạt được độ chính xác cao, đồng thời giảm đáng kể chi phí tính toán trong giai đoạn huấn luyện. Trên thực tế, tốc độ tính toán còn phụ thuộc nhiều vào mức độ tối ưu hóa phần mềm cũng như việc khai thác CUDA, vốn chưa được tối ưu hoàn toàn trong phiên bản hiện tại của thư viện FastKAN.